% =============================================================
% Fallstudien 2 - UEFA Euro - XGBoost results record
% Compile with pdflatex
% =============================================================

\documentclass[11pt,a4paper]{article}

\usepackage[margin=2.5cm]{geometry}
\usepackage{booktabs}
\usepackage{amsmath}
\usepackage{hyperref}

\title{XGBoost on the UEFA Euro data\\[0.3em]
       \large Results record -- Sheet 3, Section 1}
\author{Rohit}
\date{\today}

\begin{document}
\maketitle

\section{Setup}

\textbf{Data.} Six tournaments (2004--2024), 492 rows, 246 matches.
Each match appears twice, once from each team's perspective. The response
\texttt{Goals} has mean $1.205$ and variance $1.253$, so the Poisson
assumption holds well. All numeric predictors are differences
(team $-$ opponent).

\textbf{Objective.} \texttt{count:poisson}, tuned by Poisson negative
log-likelihood.

\textbf{Folds.} Leave-one-season-out (LOSO): one fold per tournament.
A random split would separate the two rows of a match, and the predictors
are built from pre-tournament information, so the tournament is the
natural unit.

\textbf{Match key.} \texttt{Year + sorted team pair + GroupStage}. The
\texttt{GroupStage} component is required: three team pairs met twice in
the same tournament (group stage and knockout round).

\subsection*{Hyperparameter grid}

\begin{center}
\begin{tabular}{lll}
  \toprule
  parameter & values & source \\
  \midrule
  \texttt{max\_depth}        & 1, 2, 3          & ESL ch.\ 10: 4--8 terminal nodes \\
  \texttt{eta}               & 0.01, 0.05, 0.1  & Friedman (2001): small shrinkage \\
  \texttt{subsample}         & 0.5, 0.8         & Friedman (2002) \\
  \texttt{min\_child\_weight} & 1, 5, 10         & xgboost docs \\
  \midrule
  \multicolumn{3}{l}{$3 \times 3 \times 2 \times 3 = 54$ combinations;
    \texttt{colsample\_bytree} fixed at 0.8} \\
  \bottomrule
\end{tabular}
\end{center}

\texttt{nrounds} is not in the grid. It is coupled to \texttt{eta}
(a small step size needs more rounds), so early stopping
(\texttt{early\_stopping\_rounds} $= 50$, ceiling 2000) selects it per
combination instead.

\textbf{Fits.} Nested CV: $6 \times (54 \times 5 + 1) = 1\,626$.
Final pass: $54 \times 6 + 1 = 325$.


\section{Nested cross-validation}

Outer loop over the 6 seasons (evaluation), inner LOSO over the
remaining 5 (tuning). Each outer fold re-tunes from scratch, so the
held-out season never influences the hyperparameters used to predict it.

\begin{center}
\begin{tabular}{lrrrrcccr}
  \toprule
  season & $n$ & MAE & RMSE & NLL & depth & $\eta$ & subsample & rounds \\
  \midrule
  2004 &  62 & 0.844 & 1.052 & 1.371 & 1 & 0.10 & 0.8 &  81 \\
  2008 &  62 & 0.883 & 1.078 & 1.414 & 1 & 0.10 & 0.8 &  81 \\
  2012 &  62 & 0.782 & 1.009 & 1.333 & 1 & 0.05 & 0.8 & 169 \\
  2016 & 102 & 0.846 & 1.015 & 1.313 & 1 & 0.05 & 0.5 & 172 \\
  2020 & 102 & 0.886 & 1.110 & 1.417 & 1 & 0.05 & 0.8 & 198 \\
  2024 & 102 & 0.751 & 0.968 & 1.307 & 1 & 0.05 & 0.8 & 165 \\
  \midrule
  mean &     & 0.8318 & 1.0388 & 1.3593 & & & & \\
  sd   &     & 0.0545 & 0.0515 & 0.0490 & & & & \\
  pooled &   & 0.8308 & 1.0382 & 1.3560 & & & & \\
  \bottomrule
\end{tabular}
\end{center}

All six folds chose \texttt{max\_depth} $= 1$. The selected model is
therefore purely additive: a sum of single-split stumps with no
interactions. The variation between folds is mainly in
$\eta$ and \texttt{nrounds}, which move together
($\eta = 0.1 \Rightarrow 81$ rounds, $\eta = 0.05 \Rightarrow 165$--$198$):
roughly the same total correction reached at different speeds.

2024 is the easiest season (MAE $0.751$ against a mean of $0.832$).
Training size is constant at 410 rows in every fold, so this is a
property of the tournament, not of the training data available.


\section{Selection bias}

Both columns train on five seasons and predict the sixth, on the same
folds. They differ only in whether the predicted season also helped
choose the hyperparameters.

\begin{center}
\begin{tabular}{lrrrr}
  \toprule
  metric & nested (pooled) & tuning pass (pooled) & gap & in sd \\
  \midrule
  MAE         & 0.8308 & 0.8283 & 0.0024 & 0.045 \\
  RMSE        & 1.0382 & 1.0354 & 0.0028 & 0.054 \\
  Poisson NLL & 1.3560 & 1.3533 & 0.0027 & 0.055 \\
  \bottomrule
\end{tabular}
\end{center}

All three metrics agree on a bias of roughly $0.0025$, about 4--5\% of
one standard deviation. It is small because the tuning is stable: with
depth 1 chosen in every fold, the winner barely depends on which seasons
were used to pick it. The nested figures are the ones comparable across
methods, since the tuning score flatters whoever searches the larger grid.


\section{Constant baselines}

\begin{center}
\begin{tabular}{lrr}
  \toprule
   & MAE & RMSE \\
  \midrule
  XGBoost              & 0.831 & 1.038 \\
  constant median (1)  & 0.809 & 1.114 \\
  constant mean (1.205)& 0.873 & 1.099 \\
  \bottomrule
\end{tabular}
\end{center}

MAE is minimised by the median, RMSE by the mean. The goal distribution
has median 1 and mean 1.205, so the two measures reward different
predictions. The constant median wins on MAE and loses on RMSE; the
constant mean does the reverse. A Poisson model estimates the mean, so
MAE structurally handicaps it: this is loss-function geometry, not
evidence about the predictors. Season by season, XGBoost beats the
constant mean on RMSE in all six tournaments.


\section{Final model}

Tuned by LOSO over all six seasons ($54 \times 6 = 324$ fits), then
fitted on all 492 rows.

\begin{center}
\begin{tabular}{lr}
  \toprule
  \texttt{max\_depth}         & 1 \\
  \texttt{eta}                & 0.1 \\
  \texttt{subsample}          & 0.8 \\
  \texttt{min\_child\_weight} & 5 \\
  \texttt{nrounds}            & 86 \\
  \bottomrule
\end{tabular}
\end{center}

This model carries no error of its own: it has seen every season, so any
score computed on those rows would be optimistic. The performance claim
is the nested CV estimate in Section 2.

\begin{center}
\begin{tabular}{lrr}
  \toprule
  predictor & Gain & stumps \\
  \midrule
  MarketValue & 0.392 & 26 \\
  CLPlayers   & 0.212 & 17 \\
  GDP         & 0.192 & 23 \\
  FifaRank    & 0.109 &  9 \\
  UefaPoints  & 0.096 & 11 \\
  GroupStage  & ---   &  0 \\
  \bottomrule
\end{tabular}
\end{center}

\texttt{GroupStage} is never split on, consistent with its correlation of
$-0.016$ with \texttt{Goals}. \texttt{GDP} is split on often but
contributes little per split, so frequency and importance diverge; Gain
is the column to report. The $39\%$ Gain for \texttt{MarketValue} closely
reproduces the Week 2 CART result ($\approx 40\%$).


\section{Match outcomes via the Skellam distribution}

The two predicted rates of a match are paired and converted to
win/draw/loss probabilities. With
$G_1 \sim \text{Pois}(\lambda_1)$, $G_2 \sim \text{Pois}(\lambda_2)$
independent, the goal difference $D = G_1 - G_2$ follows a Skellam
distribution:
\[
  P(\text{win}) = P(D \geq 1), \quad
  P(\text{draw}) = P(D = 0), \quad
  P(\text{loss}) = P(D \leq -1).
\]
Predictions come from the selected hyperparameters refitted per fold
(goal-level MAE $0.8296$, RMSE $1.0380$), so every match is predicted by
a model that did not see its tournament.

\begin{center}
\begin{tabular}{lrrrr}
  \toprule
   & Log loss & Brier & RPS & Accuracy \\
  \midrule
  XGBoost $+$ Skellam  & 0.9973 & 0.5982 & 0.1951 & 0.5081 \\
  Historical frequency & 1.1027 & 0.6695 & 0.2306 & 0.3008 \\
  \midrule
  improvement & 9.6\% & 10.6\% & 15.4\% & $+20.7$pp \\
  \bottomrule
\end{tabular}
\end{center}

\begin{center}
\begin{tabular}{lrrrr}
  \toprule
  tournament & Log loss & Brier & RPS & Accuracy \\
  \midrule
  2004 & 1.0116 & 0.6075 & 0.1870 & 0.4516 \\
  2008 & 1.0141 & 0.6064 & 0.2298 & 0.4516 \\
  2012 & 0.9394 & 0.5525 & 0.1905 & 0.6129 \\
  2016 & 1.0449 & 0.6296 & 0.2059 & 0.4706 \\
  2020 & 0.9661 & 0.5809 & 0.1792 & 0.5686 \\
  2024 & 0.9971 & 0.6014 & 0.1870 & 0.4902 \\
  \bottomrule
\end{tabular}
\end{center}

\subsection*{Discussion}

\textbf{Advantages.} The outcome response is what matters in football,
and it uses the whole predicted distribution rather than a point estimate.
The improvement over the baseline is far larger here (10--15\%) than at
the goal level, where XGBoost beat the ``do nothing'' floor by only about
5\% and lost to the constant median on MAE outright. The predictors carry
real information about \emph{who wins} even though they cannot pin down
\emph{how many goals}. Outcome probabilities are also directly comparable
with bookmaker odds.

\textbf{Disadvantages.} Skellam assumes the two teams' goals are
independent, which they are not: a team leading changes how it plays. The
known consequence is that draws are underestimated, and that appears here
--- mean predicted $P(\text{draw}) = 0.258$ against an actual frequency of
$0.293$, with wins correspondingly overestimated ($0.394$ vs $0.358$).
Information is also discarded: a $1$--$0$ and a $5$--$0$ become the same
``win''. The sample halves from 492 rows to 246 matches. Knockout draws
are decided on penalties, which are close to coin flips.

\textbf{Measures.} RPS is the most appropriate, because the outcomes are
ordered: predicting a loss when the result was a draw should be penalised
less than when the result was a win, and only RPS respects that. It is
also where the largest improvement appears. Log loss and Brier are proper
scoring rules and are reported alongside. Accuracy is the weakest measure
here: under Skellam a draw is never the single most likely outcome, so the
model never predicts one, and all 72 actual draws are counted wrong by
construction:

\begin{center}
\begin{tabular}{lrrr}
  \toprule
  predicted $\backslash$ actual & Win & Draw & Loss \\
  \midrule
  Win  & 67 & 43 & 28 \\
  Draw &  0 &  0 &  0 \\
  Loss & 21 & 29 & 58 \\
  \bottomrule
\end{tabular}
\end{center}

This is not the same as the model assigning draws low probability --- it
gives them $0.258$ on average. A draw simply sits between the other two
outcomes and is squeezed from both sides. The baseline's accuracy of
$0.3008$ is likewise near-meaningless: it assigns identical probabilities
to every match and so predicts a single class throughout.


\section{Notes for the group}

\begin{itemize}
  \item Comparisons require identical folds, identical metric
        definitions, and the same aggregation (pooled vs mean) across all
        four methods.
  \item The fold-to-fold standard deviation ($\approx 0.05$) exceeds the
        likely gaps between methods, so differences of a few hundredths
        should be checked with a paired test across the six folds rather
        than by comparing averages.
  \item Earlier single-split benchmarks (Lasso $0.735$, CART $0.741$ on
        2024 only) are not comparable to six-fold figures. The XGBoost
        2024 fold alone gives $0.751$ against a six-fold mean of $0.832$,
        so a 2024-only estimate is optimistic by roughly $0.08$ --- an
        order of magnitude larger than the selection bias in Section 3.
\end{itemize}

\end{document}